\documentclass[a4paper,12pt]{article}
\addtolength{\topmargin}{1cm} \addtolength{\textheight}{1cm}
\addtolength{\textwidth}{1.5cm}

% Utilizacao maior da folha
\usepackage{a4wide}
% Para escrever em Portugues
\usepackage[brazil]{babel}
\usepackage[utf8]{inputenc}
\usepackage[T1]{fontenc}

\begin{document}

\large
\begin{flushleft}
	UEL -- CCE -- Departamento de Física\\
	Disciplina: {\em MECÂNICA QUÂNTICA II} \\
	Per\'{\i}odo: $2^{\rm \underline o}$ semestre de 2021\\
	Professor: Sandro Vitenti: vitenti@uel.br, sala 7
\end{flushleft}

\normalsize

\begin{center}
	\section*{PLANO DE CURSO}
\end{center}
\section{Ementa}

A relatividade newtoniana: referenciais inerciais, transformações de Galileu. A
relatividade einsteiniana: postulados básicos, transformações de Lorentz, referenciais
inerciais. As consequências imediatas: relatividade da simultaneidade, contração do
espaço e dilatação do tempo, efeito Doppler. A geometria do espaço-tempo: espaço de
Minkowski, quadrivetores e tensores. Dinâmica relativística: equações de movimento,
quadrivetores de energia e momento, equivalência entre massa e energia. Vetores e
tensores: grandezas covariantes e contravariantes, o tensor métrico, o tensor de
Levi-Civita, gradiente, divergente e rotacional, contrações de índices, produto escalar,
equações covariantes. Correntes e densidades: leis de conservação. Eletrodinâmica:
equações de Maxwell, tensor de energia e momento, leis de conservação, força de Lorentz,
campos de cargas aceleradas, radiação, reação radiativa. Equações dos campos
relativísticos. Referenciais não-inerciais: princípio da equivalência e algumas
consequências imediatas.


\section{Objetivos}

\begin{itemize}
	\item Apresentar os fundamentos da relatividade restrita e suas implicações físicas.
	\item Introduzir a formulação tensorial e covariante da relatividade e suas
	      aplicações.
	\item Explorar a relação entre relatividade e eletromagnetismo, culminando na
	      formulação covariante das equações de Maxwell.
\end{itemize}

\section{Conteúdo Programático}

\subsection*{Primeiro Bimestre (Semanas 1-9)}
\begin{enumerate}
	\item \textbf{Introdução à Relatividade Newtoniana}
	      \begin{itemize}
		      \item Referenciais inerciais e princípio da relatividade de Galileu
		      \item Transformações de Galileu e suas limitações
		      \item Sistemas acelerados e forças fictícias
	      \end{itemize}
	\item \textbf{Postulados da Relatividade Restrita}
	      \begin{itemize}
		      \item Experimento de Michelson-Morley e falha da hipótese do éter
		      \item Postulados de Einstein
		      \item Transformações de Lorentz e suas propriedades
	      \end{itemize}
	\item \textbf{Consequências das Transformações de Lorentz}
	      \begin{itemize}
		      \item Relatividade da simultaneidade
		      \item Dilatação do tempo e contração do espaço
		      \item Paradoxo dos gêmeos e experiências reais
	      \end{itemize}
	\item \textbf{Espaço-Tempo de Minkowski}
	      \begin{itemize}
		      \item Estrutura geométrica do espaço-tempo
		      \item Intervalos espaço-temporais e cones de luz
		      \item Causalidade relativística
	      \end{itemize}
	\item \textbf{Dinâmica Relativística I}
	      \begin{itemize}
		      \item Quadrivetores e equações de movimento
		      \item Energia e momento relativísticos
	      \end{itemize}
	\item \textbf{Dinâmica Relativística II}
	      \begin{itemize}
		      \item Equivalência entre massa e energia
		      \item Transformações de energia e momento
		      \item Decaimentos relativísticos e colisões
	      \end{itemize}
	\item \textbf{Introdução a Vetores e Tensores}
	      \begin{itemize}
		      \item Definições de vetores e tensores
		      \item Índices covariantes e contravariantes
	      \end{itemize}
	\item \textbf{Operações com Tensores}
	      \begin{itemize}
		      \item Produto escalar e métrica de Minkowski
		      \item Gradiente, divergente e rotacional na notação tensorial
	      \end{itemize}
	\item \textbf{Prova 1} Avaliação do conhecimento atual dos alunos sobre os conteúdos
	      abordados
\end{enumerate}

\subsection*{Segundo Bimestre (Semanas 10-18)}
\begin{enumerate}
	\item \textbf{Tensor Energia-Momento}
	      \begin{itemize}
		      \item Definição do tensor energia-momento
		      \item Leis de conservação em relatividade
	      \end{itemize}
	\item \textbf{O Princípio da Equivalência}
	      \begin{itemize}
		      \item Introdução a referenciais não-inerciais
		      \item Campos gravitacionais fracos e efeitos relativísticos
	      \end{itemize}
	\item \textbf{Equações de Maxwell e Covariância}
	      \begin{itemize}
		      \item Equações de Maxwell no vácuo
		      \item Forma diferencial e integral
	      \end{itemize}
	\item \textbf{Trabalho} Pesquisa sobre um tema específico e apresentação no segundo
	      bimestre
	\item \textbf{Tensor Eletromagnético e Equações Covariantes}
	      \begin{itemize}
		      \item Tensor de Maxwell e equações de campo
		      \item Potencial quadrivetorial
	      \end{itemize}
	\item \textbf{Conservação de Energia e Momento no Campo Eletromagnético}
	      \begin{itemize}
		      \item Tensor energia-momento eletromagnético
		      \item Princípio da ação mínima para eletrodinâmica
	      \end{itemize}
	\item \textbf{Cargas Aceleradas e Radiação Relativística}
	      \begin{itemize}
		      \item Campos de cargas aceleradas
		      \item Radiação e reação radiativa
	      \end{itemize}
	\item \textbf{Equações de Campo Relativísticas}
	      \begin{itemize}
		      \item Introdução a equações relativísticas de campos
		      \item Aplicações na teoria de campos e relatividade geral
	      \end{itemize}
	\item \textbf{Prova Final} Avaliação abrangente dos conteúdos do curso
\end{enumerate}

\section{Metodologia}

O curso será ministrado por meio de aulas expositivas e interativas, combinando teoria e
exemplos práticos. As aulas incluirão momentos de discussão e resolução de problemas
para reforçar a compreensão dos conceitos. Além disso, serão propostas atividades de
pesquisa e exercícios para aprofundamento. As avaliações incluirão provas escritas e um
trabalho.

\section{Avaliação}

A avaliação será feita através de três notas: duas (2) avaliações realizadas
individualmente (em sala se possível) e o trabalho. Cada avaliação equivale a um terço
da nota final.

\section{Monitoria}

\begin{itemize}
	\item S. Carroll, \textit{Spacetime and Geometry}
	\item S. Weinberg, \textit{Gravitation and Cosmology}
	\item J.D. Jackson, \textit{Classical Electrodynamics}
	\item B. Schutz, \textit{A First Course in General Relativity}
	\item H. Goldstein, \textit{Classical Mechanics}
	\item L.D. Landau \& E.M. Lifshitz, \textit{Mechanics}
\end{itemize}

\end{document}
